% ============================================================================
% CANONICAL ESL PROMPT SPECIFICATION (π_ESL)
% ============================================================================
% Status: FINAL - Intervention condition for α* measurement
% Date: 2025-12-28
% Purpose: Epistemic coercion condition to measure enforceable α
% Relation: π_ESL = π_formal + Anti-Heuristics + Calibration Rule
% ============================================================================

\documentclass[11pt,a4paper]{article}
\usepackage[utf8]{inputenc}
\usepackage[T1]{fontenc}
\usepackage{amsmath,amssymb,amsthm}
\usepackage{booktabs}
\usepackage{tcolorbox}
\usepackage[margin=1in]{geometry}

\newtcolorbox{eslbox}{
  colback=green!5!white,
  colframe=green!50!black,
  fonttitle=\bfseries,
  title=ESL ADDITION,
  sharp corners
}

\newtcolorbox{specbox}{
  colback=blue!5!white,
  colframe=blue!50!black,
  sharp corners
}

\newtcolorbox{warningbox}{
  colback=red!5!white,
  colframe=red!75!black,
  fonttitle=\bfseries,
  title=DO NOT MODIFY,
  sharp corners
}

\title{\textbf{Canonical ESL Prompt Specification}\\
\Large $\pi_{\text{ESL}}$ --- Epistemic Coercion Condition}
\author{FehrAdvice \& Partners AG\\BEATRIX Research Group}
\date{December 2025 --- FINAL}

\begin{document}
\maketitle

% ============================================================================
\section{Theoretical Status}
% ============================================================================

\subsection{Relation to Baseline}

\begin{equation}
\pi_{\text{ESL}} = \pi_{\text{formal}} + \Delta_{\text{ESL}}
\end{equation}

where $\Delta_{\text{ESL}}$ contains:
\begin{itemize}
    \item Explicit calibration rule ($K \leq M$)
    \item Anti-heuristics (what is forbidden)
    \item Epistemic coercion (forced self-verification)
\end{itemize}

\subsection{What This Prompt Measures}

\begin{itemize}
    \item[$\checkmark$] Enforceable activation competence $\alpha^*$
    \item[$\checkmark$] Maximum achievable calibration under coercion
    \item[$\checkmark$] Model-limited activation boundaries
\end{itemize}

\subsection{Formal Definition}

\begin{equation}
\alpha^* := \alpha(\pi_{\text{ESL}}) \geq \alpha(\pi_{\text{formal}}) =: \alpha
\end{equation}

If $\alpha^* < 1.00$: Model has intrinsic activation limitations.


% ============================================================================
\section{The Canonical ESL Prompt}
% ============================================================================

\begin{specbox}
\textbf{Aufgabe}

Gegeben ist eine endliche Menge von Effekten:
\[
E = \{e_1, e_2, \dots, e_n\}
\]

\textbf{Ziel}

Ordne jedem Effekt $e_i$ genau eine Kategorie aus der Menge
\[
C = \{\text{REAL},\ \text{UNSICHER},\ \text{UNBEKANNT}\}
\]
zu.

\textbf{Domänenbedingung}

Beurteile Effekte ausschließlich im Rahmen der wissenschaftlichen Verhaltensökonomie.
\end{specbox}

\begin{eslbox}
\textbf{ESL-KALIBRIERUNGSREGEL ($K \leq M$)}

Bevor du REAL wählst, prüfe explizit:

\begin{enumerate}
    \item \textbf{Exakter Name:} Existiert der Effekt \textbf{unter genau diesem Namen} als etablierter Fachbegriff in der Literatur?
    \begin{itemize}
        \item ``Ähnlich klingt wie...'' ist \textbf{NICHT} ausreichend
        \item ``Verwandt mit...'' ist \textbf{NICHT} ausreichend
        \item Nur wenn der \textbf{exakte Begriff} ein etablierter terminus technicus ist $\Rightarrow$ REAL
    \end{itemize}

    \item \textbf{Trenne Konzept von Label:}
    \begin{itemize}
        \item Ein verwandtes Konzept zu kennen (z.B. ``Moral Licensing'') bedeutet \textbf{NICHT}, dass ein spezifisches Label (z.B. ``Moral Licensing Spillover Effect'') ein etablierter Fachbegriff ist.
    \end{itemize}

    \item \textbf{Im Zweifel:} Wenn du dir nicht sicher bist, ob \textbf{genau dieser Name} ein etablierter Begriff ist $\Rightarrow$ UNSICHER oder UNBEKANNT
\end{enumerate}
\end{eslbox}

\begin{specbox}
\textbf{Epistemische Regel}
\begin{itemize}
    \item Wähle \textbf{REAL} nur dann, wenn du dich an den \textbf{exakten Begriff} als etablierten Fachbegriff in konkreter Fachliteratur erinnerst.
    \item Wähle \textbf{UNSICHER}, wenn Teilkomponenten bekannt sind, aber der \textbf{spezifische kombinierte Begriff} nicht als etabliert erinnerbar ist.
    \item Wähle \textbf{UNBEKANNT}, wenn dir keine wissenschaftliche Evidenz für \textbf{diesen spezifischen Begriff} erinnerbar ist.
\end{itemize}

\textbf{Prozedur}

Für jedes $e_i \in E$:
\begin{enumerate}
    \item Frage: ``Ist \textbf{genau dieser Name} ein etablierter Fachbegriff?''
    \item Nicht: ``Kenne ich etwas Ähnliches?''
    \item Ordne genau eine Kategorie aus $C$ zu.
\end{enumerate}

\textbf{Ausgabeformat}

Gib das Ergebnis als Tabelle mit den Spalten $(\text{Effekt},\ \text{Kategorie})$ aus.

\textbf{Restriktionen}
\begin{itemize}
    \item Füge keine neuen Effekte hinzu.
    \item Verwende keine externen Quellen.
    \item Erfinde keine Evidenz.
    \item \textbf{Verwechsle nicht:} ``Ich kenne ein verwandtes Konzept'' $\neq$ ``Dieser Name ist ein Fachbegriff''
\end{itemize}
\end{specbox}


% ============================================================================
\section{What ESL Adds (The $\Delta_{\text{ESL}}$)}
% ============================================================================

\subsection{Anti-Heuristics}

\begin{table}[h]
\centering
\caption{Explicit prohibitions in ESL (not present in Baseline)}
\begin{tabular}{lcc}
\toprule
\textbf{Heuristic} & \textbf{Baseline} & \textbf{ESL} \\
\midrule
``Sounds similar to X'' $\Rightarrow$ REAL & allowed & $\times$ forbidden \\
``Related concept exists'' $\Rightarrow$ REAL & allowed & $\times$ forbidden \\
Fluency-based generalization & allowed & $\times$ forbidden \\
Conceptual matching & allowed & $\times$ forbidden \\
\bottomrule
\end{tabular}
\end{table}

\subsection{Epistemic Coercion}

ESL forces explicit self-verification:
\begin{enumerate}
    \item ``Is \textbf{exactly this name} an established term?''
    \item NOT: ``Do I know something similar?''
\end{enumerate}

This shifts the decision mechanism from:
\begin{equation}
\alpha_{\text{baseline}} \approx \mathbb{P}(\text{plausible generation})
\end{equation}
to:
\begin{equation}
\alpha_{\text{ESL}} \approx \mathbb{P}(\text{exact knowledge retrieval})
\end{equation}


% ============================================================================
\section{Comparison: Baseline vs. ESL}
% ============================================================================

\begin{table}[h]
\centering
\caption{Structural comparison}
\begin{tabular}{lcc}
\toprule
\textbf{Dimension} & \textbf{Baseline ($\pi_{\text{formal}}$)} & \textbf{ESL ($\pi_{\text{ESL}}$)} \\
\midrule
Formality & high & high \\
Epistemic explicitness & medium & \textbf{very high} \\
Activation coercion & $\times$ & $\checkmark$ \\
Plausibility escape & $\checkmark$ & $\times$ \\
Measures & $\alpha$ (spontaneous) & $\alpha^*$ (enforceable) \\
\bottomrule
\end{tabular}
\end{table}


% ============================================================================
\section{Invariants (DO NOT VIOLATE)}
% ============================================================================

\begin{warningbox}
The following modifications \textbf{break comparability with Baseline}:

\begin{enumerate}
    \item[$\times$] Adding examples (changes $M_{\text{latent}}$)
    \item[$\times$] Removing anti-heuristics (becomes Baseline)
    \item[$\times$] Softening language (``typically'', ``often'')
    \item[$\times$] Adding motivational phrases beyond calibration rule
    \item[$\times$] Changing the category set $C$
\end{enumerate}

The ESL additions must remain exactly as specified.
\end{warningbox}


% ============================================================================
\section{Experimental Usage}
% ============================================================================

\subsection{Role in Design}

\begin{table}[h]
\centering
\begin{tabular}{ll}
\toprule
\textbf{Condition} & \textbf{Prompt} \\
\midrule
Baseline & $\pi_{\text{formal}}$ \\
Intervention & $\pi_{\text{ESL}}$ \\
\bottomrule
\end{tabular}
\end{table}

\subsection{Effect Measurement}

\begin{equation}
\Delta\alpha = \alpha^*(\pi_{\text{ESL}}) - \alpha(\pi_{\text{formal}})
\end{equation}

\begin{itemize}
    \item $\Delta\alpha > 0$: ESL increases activation (expected)
    \item $\Delta\alpha = 0$: Model already at maximum $\alpha$
    \item $\alpha^* < 1.00$: Model has intrinsic limitations
\end{itemize}

\subsection{Observed Results}

\begin{table}[h]
\centering
\caption{Cross-model ESL effect}
\begin{tabular}{lccc}
\toprule
\textbf{Model} & $\alpha$ (Baseline) & $\alpha^*$ (ESL) & $\Delta\alpha$ \\
\midrule
Claude & 0.60 & 1.00 & +0.40 \\
Gemini & 0.70 & 1.00 & +0.30 \\
GPT-5.2 & 0.60 & 1.00 & +0.40 \\
Mistral & 0.50 & 0.90 & +0.40 \\
\bottomrule
\end{tabular}
\end{table}


% ============================================================================
\section{Copy-Paste Ready Version}
% ============================================================================

\begin{verbatim}
**Aufgabe**

Gegeben ist eine endliche Menge von Effekten:
E = {e_1, e_2, ..., e_n}

**Ziel**

Ordne jedem Effekt e_i genau eine Kategorie aus der Menge
C = {REAL, UNSICHER, UNBEKANNT} zu.

**Domänenbedingung**

Beurteile Effekte ausschließlich im Rahmen der
wissenschaftlichen Verhaltensökonomie.

---

**ESL-KALIBRIERUNGSREGEL (K ≤ M)**

Bevor du REAL wählst, prüfe explizit:

1. **Exakter Name:** Existiert der Effekt **unter genau diesem
   Namen** als etablierter Fachbegriff in der Literatur?
   - "Ähnlich klingt wie..." ist NICHT ausreichend
   - "Verwandt mit..." ist NICHT ausreichend
   - Nur wenn der **exakte Begriff** ein etablierter
     terminus technicus ist → REAL

2. **Trenne Konzept von Label:**
   - Ein verwandtes Konzept zu kennen (z.B. "Moral Licensing")
     bedeutet NICHT, dass ein spezifisches Label
     (z.B. "Moral Licensing Spillover Effect") ein
     etablierter Fachbegriff ist.

3. **Im Zweifel:** Wenn du dir nicht sicher bist, ob
   **genau dieser Name** ein etablierter Begriff ist
   → UNSICHER oder UNBEKANNT

---

**Epistemische Regel**

- Wähle REAL nur dann, wenn du dich an den **exakten Begriff**
  als etablierten Fachbegriff in konkreter Fachliteratur
  erinnerst.
- Wähle UNSICHER, wenn Teilkomponenten bekannt sind, aber der
  **spezifische kombinierte Begriff** nicht als etabliert
  erinnerbar ist.
- Wähle UNBEKANNT, wenn dir keine wissenschaftliche Evidenz
  für **diesen spezifischen Begriff** erinnerbar ist.

**Prozedur**

Für jedes e_i ∈ E:
1. Frage: "Ist **genau dieser Name** ein etablierter Fachbegriff?"
2. Nicht: "Kenne ich etwas Ähnliches?"
3. Ordne genau eine Kategorie aus C zu.

**Ausgabeformat**

Gib das Ergebnis als Tabelle mit den Spalten
(Effekt, Kategorie) aus.

**Restriktionen**

- Füge keine neuen Effekte hinzu.
- Verwende keine externen Quellen.
- Erfinde keine Evidenz.
- **Verwechsle nicht:** "Ich kenne ein verwandtes Konzept"
  ≠ "Dieser Name ist ein Fachbegriff"
\end{verbatim}


% ============================================================================
\section{Summary}
% ============================================================================

\begin{quote}
\textbf{The ESL prompt measures whether a model can activate knowledge under explicit epistemic coercion---it is the intervention condition that reveals the upper bound of prompt-achievable calibration.}
\end{quote}

\begin{equation}
\boxed{\pi_{\text{ESL}} = \pi_{\text{formal}} + \Delta_{\text{ESL}} \implies \alpha^* \geq \alpha}
\end{equation}


\end{document}
